\subsection{A coefficient cover and its raw-height budgets}
\label{sec:coefficient-raw-budgets}

The coefficient alternatives must be fixed before choosing a common
height cap. Their section constructions then give a raw-height cap
independently of the remaining coefficient radii. In all phase
derivatives below, the coefficients and actual finite multiplier are
held fixed.

\begin{lemma}[Coefficient-only cover]\label{lem:pure-coefficient-cover}
Fix the original nested sections, the passive compact coefficient block,
and the finite orders required by the regional preparations. There are
fixed thresholds
\(P_D,p_*,P_*,k_\flat,\Pi_{\rm far},\Pi_{\max},
\mathfrak q_\flat,U_d,U_{\rm cub},G_{\max}\), and fixed overlapping
coefficient rectangles, with the following properties.

For \(\alpha_0\ge0\), the determinant and signed-trace substitutions
\eqref{deg:determinant-chart}--\eqref{deg:trace-chart} cover every
nonzero coefficient pair. The double-zero pair is retained separately.
For \(\alpha_0=-s^2<0\), put \(p=\alpha_1/s\). The ranges
\begin{equation}\label{eq:pure-outer-coefficient-cover}
 p\le0,\qquad 0<p\le p_*,\qquad
 p_*/2\le p\le2P_*,\qquad p\ge P_*
\end{equation}
cover every coefficient fibre. The small-positive range is covered by the overlapping gamma,
compact-cubic, degree-loss and remote-root alternatives in
Table~\ref{tab:coefficient-phase-cover}.
Here the variables are exactly \eqref{deg:gamma-chart} and
\eqref{deg:cubic-chart}. The degree and remote-root substitutions are
\eqref{deg:gamma-degree-dictionary},
\eqref{deg:negative-far-dictionary}, and
\eqref{deg:cubic-degree-dictionary}. Their full independent rectangles
lie in the prescribed coefficient ranges of their preparations.
All thresholds and reserves are chosen before a final raw coefficient
restriction. Neither a phase input, a return zero, nor a height cap is
used to choose a coefficient alternative.
\end{lemma}

\begin{proof}
Fix a signed compact strict source block about \(a=p=0,Y(0)=1\).
Choose its high-height constant \(K_{\rm hi}\), large enough also for
the signed trace block \(1/2\le|p_L|\le3/2\), and put
\(K_{\rm lo}=2K_{\rm hi}\). The trace preparation on an initial
\(|a_L|\le a_0\le1\) uses the fixed central-height range
\(K_{\rm lo}(1+(3/2)^2)\). Its graph barrier fixes the bounded
actual-gap range; its terminal construction then fixes
\(Q,R_0,a_*,\xi_0\). Choose \(P_D>2/\sqrt{a_*}\).
When \(\alpha_0>0\), use \(r_D=\sqrt{\alpha_0}\) if
\(|\alpha_1|\le P_D\sqrt{\alpha_0}\). Otherwise put
\[
 r_L=|\alpha_1|,\qquad p_L=\operatorname{sgn}\alpha_1,
 \qquad a_L=\alpha_0/\alpha_1^2<a_*/4.
\]
The two coefficient descriptions overlap on
\(P_D/2<|\alpha_1|/\sqrt{\alpha_0}<P_D\).
For zero opening and nonzero trace, the second description has
\(a_L=0\).

For negative opening, first choose the bounded positive-quadratic
root-free and ghost reserves, and the independent positive far-root
bounds \(q_0,\lambda_0,\tau_0\) in
\eqref{deg:positive-far-range}. Choose the reciprocal preparation
bounds \(\mathfrak q_-,\Pi_-,\tau_-\), with its positive-unit lower
bound \(\omega_{\min}\), as in
\eqref{deg:negative-far-rectangle}. Put
\begin{equation}\label{eq:pure-routing-reserves}
\begin{gathered}
\mathfrak q_\flat=\tfrac12\min\{q_0,\mathfrak q_-\},\qquad
\Pi_{\rm far}=\max\{\lambda_0^{-1/2},\Pi_-,2\sqrt{12}\},\\
\Pi_{\max}=2\Pi_{\rm far}.
\end{gathered}
\end{equation}
On \(0<\pi\le\Pi_{\max}\), fix the two bounded-gamma preparations,
including the positive ghost raw-tail reserve. Choose \(k_\flat>0\)
so that \(2k_\flat\) lies below their fixed \(k\)-bounds and below
\(k_d(\mathfrak q_\flat)\). The latter is the bound placing the
gamma-degree dictionary in its prescribed compact positive-unit block
for \(\mathfrak q\ge\mathfrak q_\flat\). Set
\(G_{\max}=k_\flat^{-2/3}\). Fix \(U_d\) for the large-\(U\)
cubic-degree dictionary on this \(G\)-block and take
\(U_{\rm cub}>U_d+1\). Fix the compact cubic preparation on
\([0,U_{\rm cub}]\times[-G_{\max},G_{\max}]\).

These substitutions are algebraic. The gamma polynomial
\(t^3+3\epsilon t^2-\mathfrak qt+k^2\) and Vieta's relations give
\eqref{deg:gamma-degree-dictionary}; the positive-unit entries have
compact closure for \(\mathfrak q\ge\mathfrak q_\flat\).
On the negative far face, the same relations give
\(\omega=1-a+ja\), \(\tau=p^2/\omega^2\), and the full rectangle
\eqref{deg:negative-far-rectangle}. The large-\(U\) polynomial
\(y^3+3Gy^2-Uy+1\) gives
\eqref{deg:cubic-degree-dictionary}; its positive units tend uniformly
to one for bounded \(G\). Let \(\tau_d>0\) be the resulting compact
degree bound. Take both degree radial bounds below \(1/2\), including
their bounded-source, enlarged positive-energy, raw-image and relative
curvature reserves. Include the fixed core/high overlap reserve in
\(\tau_0\) and choose once
\begin{equation}\label{eq:pure-small-ratio-reserve}
0<p_*^2<\tfrac12\min\{\tau_0,\tau_d,
                              \tau_-\omega_{\min}^2\}.
\end{equation}
Thus \(\tau=p^2\) and \(\tau=p^2/\omega^2\) already lie in their
respective independent rectangles.

If \(\gamma\ne0\) and \(k\le2k_\flat\), the overlapping
\(\pi\)-ranges are \(\pi\le\Pi_{\max}\) and
\(\pi\ge\Pi_{\rm far}\). On the latter range split at
\(\mathfrak q_\flat\). The large-\(\mathfrak q\) dictionary is in
the compact degree block, and the small-\(\mathfrak q\) dictionaries
are in the positive far or negative reciprocal rectangle according to
the sign of \(\gamma\). If \(\gamma=0\) or \(k\ge k_\flat\),
then \(|G|^{3/2}\le k_\flat^{-1}\), and
\([0,U_{\rm cub}]\cup[U_d,\infty)\) covers \(U\ge0\).
The \(k\), \(\pi\), \(\mathfrak q\), and \(U\) descriptions
have the strict overlaps just chosen, with their equality faces retained.

For the large positive-source range use the exact trace dictionary
\eqref{deg:large-trace-dictionary}. Its signed opening is
\(a_L=-p_L^2/p^2\). Choose
\(P_*>\max\{2,p_*,3/\sqrt{a_*}\}\), increasing it if necessary
for the fixed bounded-gap trace block of
Lemma~\ref{lem:large-ratio-height}. Finally fix the compact positive
interval \([p_*/2,2P_*]\) and its endpoint, middle and high reserves.
The remaining independent radii \(r_D,r_L,g,t,s\) can subsequently
be restricted in raw coordinates. This proves the coefficient cover.
\end{proof}

\paragraph{The small-positive-source subcover.}
\label{sec:cover-index}\label{subsec:cover-index}
The outer overlaps are \(P_D/2<|\alpha_1|/\sqrt{\alpha_0}<P_D\),
\(p_*/2<p<p_*\), and \(P_*<p<2P_*\). With
\(H=z(\gamma)\), the low and high estimates overlap on
\(K_{\rm hi}(|\alpha_0|+\alpha_1^2)<H<
K_{\rm lo}(|\alpha_0|+\alpha_1^2)\), where
\(K_{\rm lo}=2K_{\rm hi}\). Every phase description is restricted
to the common receiver of Lemma~\ref{lem:common-raw-receiver}.

In the first three rows assume \(\gamma\ne0\), \(k\le2k_\flat\),
using \eqref{deg:gamma-chart}. The next two use \eqref{deg:cubic-chart}
when \(\gamma=0\) or \(k\ge k_\flat\), with \(|G|\le G_{\max}\).
\begingroup
\let\vss\vfil
\small
\setlength{\tabcolsep}{4pt}
\renewcommand{\arraystretch}{1.08}
\begin{longtable}{@{}>{\raggedright\arraybackslash}p{.31\linewidth}>{\raggedright\arraybackslash}p{.39\linewidth}>{\raggedright\arraybackslash}p{.23\linewidth}@{}}
\caption{The small-positive-source coefficient subcover. Each regional estimate
counts the whole physical fibre.}\label{tab:coefficient-phase-cover}\\
\hline
Range & Regional estimate & Overlap or boundary\\
\hline
\endfirsthead
\multicolumn{3}{l}{\tablename\ \thetable\ (continued)}\\
\hline
Range & Regional estimate & Overlap or boundary\\
\hline
\endhead
\hline
\endfoot
\(0<\pi\le\Pi_{\max}\)
 & Proposition~\ref{sharp:upper-gamma}: positive gamma \(3\), negative gamma \(2\)
 & \(\Pi_{\rm far}<\pi<\Pi_{\max}\); real, fold, ghost and root-free faces\\
\(\pi>\Pi_{\rm far}\), \(\mathfrak q\ge\mathfrak q_\flat\)
 & \eqref{deg:gamma-degree-dictionary}; Proposition~\ref{p2:prop-physical-degree-count}: \(2\)
 & Far receivers extend past \(\mathfrak q_\flat\)\\
\(\pi>\Pi_{\rm far}\), \(\mathfrak q<\mathfrak q_\flat\)
 & Positive gamma: \eqref{deg:far-negative-budget} and Proposition~\ref{sharp:upper-gamma}; negative gamma: \eqref{deg:negative-far-dictionary} and Proposition~\ref{p2:prop-physical-degree-count}; both \(2\)
 & \(\mathfrak q_\flat<\mathfrak q_-\), \(\Pi_{\rm far}\ge\Pi_-\)\\
\(0\le U\le U_{\rm cub}\)
 & Proposition~\ref{sharp:upper-cubic}: \(3\), including the original raw ghost tail
 & \(k_\flat<k<2k_\flat\), \(U_d<U<U_{\rm cub}\)\\
\(U\ge U_d\)
 & \eqref{deg:cubic-degree-dictionary}; Proposition~\ref{p2:prop-physical-degree-count}: \(2\)
 & \(U_d<U<U_{\rm cub}\)\\
\end{longtable}
\endgroup

The degree-loss phase subcover and its internal dividers are those of
Proposition~\ref{p2:prop-unified-carrier-cover}. Each row selects a
whole-fibre bound; Proposition~\ref{sharp:upper-coefficient-bounds}
will take their maximum after the common raw restrictions.


\begin{proposition}[Regional raw-height budgets]
\label{prop:regional-raw-budgets}
Fix the coefficient rectangles and all section and overlap reserves of
Lemma~\ref{lem:pure-coefficient-cover}. Let \(\mathcal A_j\) be their
finitely many regional alternatives and write
\[
\mathcal V_\delta=\{\max(|A-1|,|C-1|,|\gamma|,
                                  |\alpha_0|,|\alpha_1|)<\delta\}.
\]
There is \(b>0\) such that, for every
\(0<h\le b\) and each \(j\), there are independent raw bounds
\(\delta_j(h)>0\) and \(\bar m_j(h)>0\) for which the actual regional
transition identities and complete derivative and curvature estimates
hold on their stated phase domains for every
\(\theta\in\mathcal A_j\cap\mathcal V_{\delta_j(h)}\) and every
\(0\le m\le\bar m_j(h)\). These domains together cover every input of
\begin{equation}\label{eq:regional-raw-budget-domain}
\{x:a_\theta+x\in J_\theta^{\rm out},\quad
                         0<h_\theta(a_\theta+x)<h\}.
\end{equation}
The bounds restrict \((A,C,\gamma,\alpha_0,\alpha_1)\) in the
independent raw coordinates and \(0\le m\le\bar m_j(h)\).
These additional restrictions impose no relation between active radii
or root gaps.
Each estimate retains its original input and inverse output windows,
positive-time passage, speed and stop clips, admitted equality dividers,
and intervening nonroot inputs. Its raw observers and divided closing
are composed on their full admitted images, with ordinary phase
\(C^3\) bounds and all nonlinear curvature terms retained.
\end{proposition}

\begin{proof}
We first construct a fixed signed high-height cap on the original raw
sections. Fix the enlarged signed strict source block used in the
coefficient lemma. Choose the core cut \(L\) and raw port radius
\(0<d<\delta\), with \(d\le1/100\), as in
Lemma~\ref{ann:p0-fixed-tail-ports}. Let \(a_H,p_H\) bound its
normalized \(|a|,|p|\), and increase \(L\) until
\(p_H/L\le1/40\), \(a_H/L^2\le1/4000\).
The fixed-core continuity bounds \(s_{\rm core},\rho_{\rm core}>0\)
put both entrance ratios in \([1/4,3/4]\) for independent
\(s\le s_{\rm core}\) and
\(|\gamma|,|A-1|,|C-1|\le\rho_{\rm core}\).

Choose a fixed transverse tube from these ports to the original
sections, and let \(\mu_z>0\) be smaller than its port tolerance and
both raw-window margins. With \(T=\log(d/(sL))\), the entrance-memory
estimate of Lemma~\ref{ann:p0-fixed-tail-ports}, with its fixed
\(\lambda=7/5\), gives
\begin{equation}\label{eq:raw-high-entrance-memory}
 |w_\sigma(T)-w_\Gamma(d)|
 \le C\{s+|\gamma|+d|A-1|+d^2|C-1|\}.
\end{equation}
Indeed its initial term is \(O(e^{-\lambda T})=O(s)\), and
integration of the two source terms has orders \(e^{-T}\) and
\(e^{-\lambda T}\); the passive terms have the three orders displayed.
The constant depends on the fixed ports and strict block, not on the
flight length. Choose \(s_H>0\) and independent passive coefficient
bounds so that the right side times \(d^2\), plus \(2d|\gamma|\),
is less than \(\mu_z\). The invariant ratio strip from the same lemma
keeps both complete tails in the chart. Their raw port errors are
therefore below \(\mu_z\), so both full images lie in the original
windows. The regular observers and divided closing can now be
composed on their larger fixed buffers.

At any signed high input freeze \(r_0=\sqrt H\). The transformed
parameters satisfy \(|a|\le K_{\rm hi}^{-1}\),
\(|p|\le K_{\rm hi}^{-1/2}\), and \(Y(0)=1\).
If \(H\le s_H^2\), this is the same admitted strict block.
The two absolute tail variations in
\eqref{sharp:upper-tail-ratio} have the common
\(r_0^{2-\lambda_\gamma}\) factor; its cancellation, followed by
the actual bridges and divided closing, gives a fixed \(C_H\).
The dilation is frozen during differentiation. This proves the signed
high estimate without a restriction on \(\gamma\log r_0\).

Let \(z_{\rm geo}>0\) be a fixed central reserve below the upper
chart boundary constructed before Lemma~\ref{lem:loc-through}, and
bound \(0<\ell_\theta(\gamma)\le L_{\rm raw}\) on the preliminary
raw block. Choose
\begin{equation}\label{eq:regional-fixed-raw-cap}
b=\tfrac14\min\{s_H^2,z_{\rm geo},d_{\rm out}/(4L_{\rm raw})\}>0.
\end{equation}
Every signed high input below \(b\) consequently has its original
raw-domain description. With \(|\alpha_0|<d_{\rm out}/4\), its
speed satisfies
\(\sigma_\theta\le|\alpha_0|+L_{\rm raw}H<d_{\rm out}/2\).

Now fix any \(0<h\le b\). The lower and strict preparations have
already fixed their compact blocks, chords and positive raw-window
reserves. We place their comparison and upper shots below \(h\)
by the following independent raw restrictions. In the determinant and
trace charts it is enough to require
\begin{equation}\label{eq:regional-det-trace-cap-bounds}
r_D^2<\frac{h}{8K_{\rm hi}(1+P_D^2)},\qquad
r_L^2<\frac{h}{4K_{\rm lo}}.
\end{equation}
Their bounded normalized-height complements have the finite strict
domains of \eqref{deg:det-height-cover} and
\eqref{deg:trace-domain-inclusion}; their larger heights are in the
signed high block. On the fixed compact positive-ratio region the
enlarged strict box \([u_L/2,3u_H]\) is placed below \(h/4\) by
\(s^2<h/(12u_H)\). The endpoint admission and its actual matching
are those of Lemma~\ref{ann:lem-compact-positive-domains}.

For a bounded-gamma preparation let \(k_B,\Pi_B\) be its fixed
coefficient bounds and \(E_B>0\) its fixed upper core height. Retain
the action comparison bound and impose
\begin{equation}\label{eq:regional-gamma-cap-bounds}
g<\min\left\{1,\frac{E_B}{4K_{\rm hi}(1+\Pi_B^2)},
                       \frac1{2k_B}\sqrt{\frac h{E_B}}\right\}.
\end{equation}
Then the normalized high threshold is below \(E_B/2\), while
\((gk)^2E_B<h/4\). The actual comparison and upper shots and every
intervening graph retain the outward admission used in
\eqref{deg:negative-gamma-cover} and in the positive paired
construction. The positive ghost raw tail is part of the same original
domain, including \(gku\ge d_0\); its raw derivative estimate is
\eqref{bp-tail:consumed-derivative}.

For the compact cubic block let \(E_c>0\) be its fixed upper core
height. Retain the exponentially small comparison bound and require
\begin{equation}\label{eq:regional-cubic-cap-bounds}
\tau_{\rm cub}<\min\left\{1,
 \frac{E_c}{4K_{\rm hi}(1+U_{\rm cub}^2)},
 \tfrac12\sqrt{h/E_c}\right\}.
\end{equation}
Here \(\tau_{\rm cub}=(-\alpha_0)^{1/3}\). Thus the normalized
threshold is below \(E_c/2\) and
\(\tau_{\rm cub}^2E_c<h/4\). The comparison remains in the actual
lip domain, with the bounded/growing resection and the original raw
ghost continuation retained in \eqref{deg:compact-physical-cover}.

For the positive far-root rectangle, the fixed \(\tau_0\) already
ensures \(\tau K_{\rm hi}(1+\tau)<E_f/2\) and the action comparison,
where \(E_f>0\) is the selected upper shot. Since
\(a=gq\lambda\), \(q\le q_0\), \(\lambda\le\lambda_0\), the
raw restriction
\begin{equation}\label{eq:regional-far-cap-bound}
g<\frac1{2q_0\lambda_0}\sqrt{h/E_f}
\end{equation}
gives \(a^2E_f<h/4\). The full lower output and upper core images
are the admitted images used in \eqref{deg:far-physical-cover}.
No ratio threshold is changed after \(h\) is fixed.

For the degree preparations, the exact source dictionary
\eqref{deg:degree-source} gives
\begin{equation}\label{eq:degree-cap-raw-identities}
\alpha_1=\tau^2\kappa^2\mathcal R\omega,\qquad
s_s^2=\tau^3\kappa^4\mathcal R^2B.
\end{equation}
The actual \(h_L=0\) source graph and its real enlarged chord have
\(Y_0\le C_Y(1-\log(2\tau))\) by
\eqref{bd:positive-height} and the enlarged positive-energy construction
in Lemma~\ref{p2:lem-bounded-source-budget}. Use the degree bounds
\(\kappa\le\kappa_*\), \(\tau\le\tau_*<1/2\) already fixed in
the coefficient cover, and put
\begin{equation}\label{eq:degree-cap-compact-constant}
D_*=\tau_*[1-\log(2\tau_*)],\qquad
C_d=C_Y\kappa_*^2
       \frac{\mathcal R_{\max}B_{\max}}{\omega_{\min}}D_*.
\end{equation}
The function \(t[1-\log(2t)]\) increases on \((0,1/2)\).
Consequently the actual trace \(h^0\) satisfies
\begin{equation}\label{eq:degree-cap-height-raw-bound}
H_\theta(h^0)\le C_Y\alpha_1\tau\kappa^2
       \frac{\mathcal RB}{\omega}[1-\log(2\tau)]
       \le C_d\alpha_1.
\end{equation}
On the reciprocal graph \eqref{deg:degree-infinity-graph}, write
\(a_{\rm rec,*}=\kappa_*A_{d,\max}<1/2\) for its fixed opening bound.
There
\[
\frac{\mathcal RB}{\omega}
=\frac{dA_dj}{1-a+ja}
\le\frac{d_{\max}A_{d,\max}j_0}{1-a_{\rm rec,*}}.
\]
Thus \eqref{eq:degree-cap-height-raw-bound} holds with
\begin{equation}\label{eq:degree-cap-reciprocal-constant}
C_r=C_Y\kappa_*^2
      \frac{d_{\max}A_{d,\max}j_0}{1-a_{\rm rec,*}}D_*
\end{equation}
in place of \(C_d\), uniformly as \(j\downarrow0\). The independent
raw restriction
\begin{equation}\label{eq:degree-cap-alpha-one-bound}
0<\alpha_1<\frac h{4\max\{C_d,C_r\}}
\end{equation}
places the upper source shot and its chosen enlarged chord below
\(h/4\). The fixed positive-label threshold also gives
\(H_\theta(h^0)>K_{\rm hi}(|\alpha_0|+\alpha_1^2)\), so that this
chord overlaps the same signed high continuation.
The fixed \(p_*\) and degree ratios are unchanged.

Raw output admission in these degree charts is independent of a
reciprocal lower bound on a gain. Indeed
\(r_s^2=(\mathcal R/\omega)\alpha_1\le C\alpha_1\); for a fixed
\(|\lambda_\gamma|<\nu<1\), the retained absolute factor obeys
\[
\Omega=r_s^{2-\lambda_\gamma}
       \le C\alpha_1^{1-\nu/2}\longrightarrow0.
\]
The stopped raw trace has a fixed window margin, so the full enlarged
output chord is admitted before its observer is composed. The exact
input buffer, actual matching and every-point allocation are those of
Proposition~\ref{p2:prop-unified-carrier-cover}. Actual reciprocal
maps use \(j>0\); \(j=0\) retains its one-sided normalized meaning.
The raw-Jacobian estimate already holds on the full original raw-window
and inverse-output-window domain and can be restricted to any such cap.

All these descriptions are restrictions of the original physical map
\(\Pi=\mathcal C_{\theta,m}\circ T_\theta(a_\theta+\cdot)\).
For local nonlinear observers write
\(\Pi=\psi_-\circ P\circ\iota\), where
\(\iota=\psi_+^{-1}\) on the actual incoming image. In particular
the physical curvature contains all three terms
\begin{equation}\label{eq:raw-budget-complete-curvature}
\Pi''=\psi_-''(P)(P')^2(\iota')^2
       +\psi_-'(P)P''(\iota')^2
       +\psi_-'(P)P'\iota''.
\end{equation}
Here the derivatives of \(P\) are evaluated at \(\iota(x)\). The regional
estimates include these ordinary observer derivatives, through
order three, before any positive affine normalization
\eqref{sharp:upper-affine}. Restriction leaves those estimates unchanged.
The divided closing has the ordinary \(C^3\) rows
\eqref{phys:eq-closing-jets}, including \(m=0\).

Take the finite minimum of the original raw-domain and error bounds
and the independent restrictions above. Express
\(\tau_{\rm cub}<T\) as \(|\alpha_0|<T^3\),
\(r_D<R\) as \(\alpha_0<R^2\), and
\(r_L<R\) as \(|\alpha_1|<R\). Half this minimum defines
\(\delta_j(h)>0\) in the independent raw coordinates. All other
normalized coordinates remain in the rectangles already fixed by
Lemma~\ref{lem:pure-coefficient-cover}.

Lastly let \(C_j\) bound the complete derivative on the regional
bounded pieces and let \(M_j\) bound the full translated output on
the original closing image. Intersect the ordinary closing-domain
bound, the regional critical-logarithm bound, and the inequalities
\(mC_j<1/2\), \(mM_j<a_{\rm cl}/2\). Their finite positive
minimum gives \(\bar m_j(h)>0\). The enlarged domain and equality
buffers were chosen earlier; their estimates hold at nonclosed inputs
as well. This proves the asserted raw budgets.
\end{proof}
